\section{Chains and the main theorem}
\label{sec:chains}

This section chooses the chain and its parameters and proves the main
theorem (Theorem~\ref{thm:main}). Next to the terminal pair, an
\emph{arity ramp} climbs the diagonal (Section~\ref{subsec:chains-ramp});
above it, a \emph{bulk} of stages of one arity \(D\) makes the source
pair grow geometrically (Section~\ref{subsec:chains-bulk}).
Section~\ref{subsec:chains-potential} shows that the chain schemes of
Section~\ref{subsec:chains-legal} have \(p>9/4\).

\subsection{Chain schemes}
\label{subsec:chains-legal}

Fix a chain (Section~\ref{subsec:chain-counts}) with source pair
\((r_1,b_1)\), evaluation size \(s=r_1+b_1-1\) and terminal pair
\((1,t)\); we call \(G=\sum_{i=1}^J\log\frac{d_i+2}{d_i}\) its
\emph{gain}. For integers \(N,k\ge1\) and, at each stage \(i\), a type
of width \(N\) with sector count \(M_i\), Proposition~\ref{prop:chain-counts} gives a
restriction \(r_*\MM_A\restr m_*\MM_B\) over \(\Q\) with the
counts~\eqref{eq:stage-chain-counts}. If \(B>A\), it is a relative
scheme, which we call a \emph{chain scheme}
(Definition~\ref{def:class} extends the term to crops, exchanges,
several blocks per stage and other supplies). Dividing the numerator and
the denominator of its pooled exponent \(p=\log(r_*/m_*)/\log(B/A)\) by
\(2Nk\) gives
\begin{equation}\label{eq:p-normalized}
 p=\frac{3\log s+C_{\mathrm{sep}}+C_{\mathrm{ex}}}{\Delta},\qquad
 \Delta=\frac{\log(B/A)}{2Nk}
 =\log t+\sum_i\Bigl(\frac{\log M_i}N-\log d_i\Bigr),
\end{equation}
where \(C_{\mathrm{sep}}=\frac3N\sum_i\log(2-1/M_i)\) is the
\emph{separator charge} and
\(C_{\mathrm{ex}}=\frac1{2Nk}\sum_i\log E_i\) is the
\emph{exactification charge}. The fixed-batch bound~\eqref{eq:prelim-phi} is
\begin{equation}\label{eq:chains-phi}
 \Phi=\frac{\log(2r_*/m_*)}{\log(B/A)}=p+\frac{C_{\mathrm{pool}}}{\Delta},
 \qquad C_{\mathrm{pool}}=\frac{\log2}{2Nk}.
\end{equation}
For large \(N\) and \(k\) the charges are small, and for types with the
largest sector counts \(\Delta\) is close to \(\log t+G\)
(Lemma~\ref{lem:types}). So the best a chain can give is its
\emph{ideal exponent} \(3\log s/(\log t+G)\);
Lemma~\ref{lem:profile-potential} shows that \(p\) always exceeds it.

\begin{lemma}[Smallest parent]\label{lem:smallest-parent}
Let \(d\ge2\), and let \((c,r)\) be a pair of positive integers. Put
\[
 b=\Bigl\lceil\frac{(d+2)c+r-1}{d}\Bigr\rceil,\qquad
 e=db-(d+2)c-r+1 .
\]
Then \(0\le e\le d-1\), the stage \((r,b)\to(c,r)\) of arity \(d\) is
legal, and \(b\) is the least integer with this property.
\end{lemma}

We call \((r,b)\) the \emph{smallest parent} of \((c,r)\) for arity
\(d\), and \(e\) its \emph{slack}.

\begin{proof}
The bounds on \(e\) restate the definition of the ceiling. For the
stage \((r,b)\to(c,r)\), the last condition of~\eqref{eq:stage-conditions}
reads \(db\ge(d+2)c+r-1\), that is, \(e\ge0\), and the ceiling is the
least integer \(b\) with this property.
\end{proof}

\subsection{The ramp}
\label{subsec:chains-ramp}

Recall that the least profile of~\citep[(3.4)]{FinCon26} is
\(F(m,n)=\varsigma_{\min(m,n)}\bigl(\max(m,n)+\frac{\min(m,n)-1}2\bigr)\),
where \(\varsigma_m=\prod_{\ell=1}^{m-1}(1+\frac1{3\ell})\) (written \(s_m\)
in~\citep{FinCon26}), so that \(F(R,R)=\varsigma_R(3R-1)/2\).

\begin{lemma}[Zero-slack ramp]\label{lem:ramp}
Let \(R\ge2\). Start at the pair \((1,1)\) and, for
\(j=1,\ldots,R-1\) in turn, take the smallest parent for arity \(3j-1\)
and then the smallest parent for arity \(3j\). All slacks are zero, and
the pairs are
\[
 (j,j)\xleftarrow{\;3j-1\;}(j,j+1)\xleftarrow{\;3j\;}(j+1,j+1)
 \qquad(1\le j\le R-1),
\]
where each arrow points from a parent to its child and carries the
arity. Read forwards, this is a chain of \(2(R-1)\) legal stages from
\((R,R)\) to the terminal pair \((1,1)\), with largest arity \(3R-3\).
Its gain is
\[
 G=\sum_{j=1}^{R-1}\log\frac{(3j+1)(3j+2)}{(3j-1)\,3j}=\log F(R,R),
 \qquad\text{and}\qquad F(R,R)\ge R^{4/3}.
\]
\end{lemma}

\begin{proof}
For arity \(3j-1\) and the pair \((j,j)\),
Lemma~\ref{lem:smallest-parent} gives
\(b=\lceil((3j+1)j+j-1)/(3j-1)\rceil=\lceil(3j-1)(j+1)/(3j-1)\rceil=j+1\)
with slack zero. For arity \(3j\) and the pair \((j,j+1)\), it gives
\(b=\lceil((3j+2)j+j)/(3j)\rceil=j+1\) with slack zero. The two stages
contribute the factors \((3j+1)/(3j-1)\) and \((3j+2)/(3j)\) to
\(\exp(G)\). Since \(3j+2=3(j+1)-1\), the product of the factors
\((3j+2)/(3j-1)\) over \(j<R\) telescopes to \((3R-1)/2\), and the
product of the factors \((3j+1)/(3j)\) is \(\varsigma_R\). Hence
\(\exp(G)=F(R,R)\). Finally, \((1+\frac1{3\ell})^3\ge1+\frac1\ell\)
gives \(\varsigma_R^3\ge\prod_{\ell=1}^{R-1}(1+\frac1\ell)=R\), and
\((3R-1)/2\ge R\).
\end{proof}

Each pair of ramp stages multiplies \(\exp(G)\) by
\(F(j+1,j+1)/F(j,j)\), and \(F(j+1,j+1)-F(j,j)=\varsigma_j+\varsigma_{j+1}\).
For \(F\), the quantity \(H_a\) in the proof
of~\citep[Lemma~5.1]{OAI26} equals \(\varsigma_a\), so the ramp attains the
diagonal recursion~\citep[(5.1)]{OAI26} with equality, with integer
pairs and no rounding. Used alone, the ramp is a chain with \(s=2R-1\),
and its ideal exponent \(3\log(2R-1)/\log F(R,R)\) is the profile
exponent of~\citep[Section~3]{FinCon26}. It exceeds \(9/4\) by less
than \(\frac94\log2/\log R\), since \(F(R,R)\ge R^{4/3}\)
(Lemma~\ref{lem:ramp}), and, as the proof of
\citep[Proposition~11.1(b)]{FinCon26} shows, by at least
\(\frac94\log\frac65/\log\frac{3R-1}2\). So, by
Lemma~\ref{lem:profile-potential}(c), a chain scheme on the ramp alone
with \(p\le9/4+\eps\) needs \(R\) and the number of stages exponential
in \(1/\eps\). The bulk is what makes \(\log A\) polynomial in
\(1/\eps\).

\subsection{The bulk}
\label{subsec:chains-bulk}

\begin{lemma}[Bulk potential]\label{lem:bulk-potential}
Let \(D\ge2\), let \(\mu\) be the positive root of
\(D\mu^2-\mu-(D+2)=0\), and put \(\eta=\mu-1/D\). Then
\(1+1/D<\mu<5/4+1/D\), so \(1<\eta<5/4\), and \(\mu\eta=(D+2)/D\). For a
pair \((r,b)\) put \(V(r,b)=\eta(r-\frac13)+(b-\frac13)\).
\begin{enumerate}
\item[(a)] If \((r,b)\) is the smallest parent of \((c,r)\) for arity
\(D\), with slack \(e\), then \(V(r,b)=\mu V(c,r)+e/D\).
\item[(b)] Let \(R\ge1\) and \(L\ge0\), and let the pair \((r',b')\)
arise from \((R,R)\) by \(L\) smallest-parent steps of arity \(D\). Then
\(r'+b'-1<\mu^L(\frac94R+D)\).
\end{enumerate}
\end{lemma}

\begin{proof}
Put \(f(u)=Du^2-u-(D+2)\). Since \(f(0)<0\) and \(f\) is a convex
quadratic, \(f\) has exactly one positive root,
\(\mu=(1+\sqrt{4D^2+8D+1})/(2D)\). As \(f(1+\frac1D)=-1\) and
\(f(\frac54+\frac1D)=(9D-12)/16>0\), it lies strictly between
\(1+\frac1D\) and \(\frac54+\frac1D\). The equation \(D\mu^2=\mu+D+2\)
gives \(\mu\eta=\mu^2-\mu/D=(D+2)/D\).

(a) By the definition of the slack, \(b=((D+2)c+r-1+e)/D\). In
\[
 V(r,b)-\mu V(c,r)
 =\eta\bigl(r-\tfrac13\bigr)+\frac{(D+2)c+r-1+e}D-\tfrac13
 -\mu\eta\bigl(c-\tfrac13\bigr)-\mu\bigl(r-\tfrac13\bigr),
\]
the coefficient of \(c\) is \((D+2)/D-\mu\eta=0\), and that of \(r\) is
\(\eta+1/D-\mu=0\). The constant term is
\(\frac{\mu-\eta}3+\frac{\mu\eta-1}3+\frac{e-1}D
=\frac1{3D}+\frac2{3D}+\frac{e-1}D=\frac eD\).

(b) Let \(V_\ell\) be the value of \(V\) after \(\ell\) steps, and
\(e_\ell\le D-1\) the slack of step \(\ell\). By (a),
\[
 V_L=\mu^LV_0+\sum_{\ell=1}^L\mu^{L-\ell}\frac{e_\ell}D
 \le\mu^LV_0+\sum_{\ell=0}^{L-1}\mu^\ell
 \le\mu^L\Bigl(V_0+\frac1{\mu-1}\Bigr).
\]
Here \(V_0=(1+\eta)(R-\frac13)<\frac94R\), and \(\mu-1>1/D\) gives
\(1/(\mu-1)<D\). Finally, \(\eta>1\) and \(r'\ge1\) give
\(V_L=\eta(r'-\frac13)+b'-\frac13>r'+b'-1\).
\end{proof}

The potential \(V\) comes from the map \((c,r)\mapsto(r,((D+2)c+r-1)/D)\),
the smallest parent without the ceiling: the map fixes
\((\frac13,\frac13)\), and \((\eta,1)\) is a left eigenvector of its
linear part, for the Perron root \(\mu\).

Each bulk stage adds \(g_D=\log\frac{D+2}D\) to the gain and multiplies
\(s\) by about \(\mu\): in the setting of
Lemma~\ref{lem:bulk-potential}(b), part (a) and \(e\ge0\) also give
\(V(r',b')\ge\mu^LV(R,R)\), and \(V(r,b)<\frac54r+b\le\frac52(r+b-1)\)
for every pair of positive integers, so
\(\log(r'+b'-1)=L\log\mu+O(\log(R+D))\). For a long bulk,
\(3\log s/G\) therefore tends to \(\omega_D=3\log\mu/g_D\), the ideal
exponent of the bulk.

\begin{lemma}[Ideal gap]\label{lem:ideal-gap}
Let \(D\ge2\), let \(\mu\) be as in Lemma~\ref{lem:bulk-potential}, and
put \(y=g_D/2\). Then \(\sinh y=1/\sqrt{D(D+2)}\),
\[
 \log\mu=y+\operatorname{arsinh}\bigl(\tfrac12\sinh y\bigr),
 \qquad\text{and}\qquad
 \frac94\le\omega_D\le\frac94+\frac{27}{32(D+1)^2}.
\]
\end{lemma}

\begin{proof}
Since \(e^y=\sqrt{(D+2)/D}\), we have
\(\sinh y=\frac12(e^y-e^{-y})=1/\sqrt{D(D+2)}\). With
\(z=\frac12\sinh y\),
\[
 z+\sqrt{1+z^2}=\frac{1+\sqrt{4D^2+8D+1}}{2\sqrt{D(D+2)}},
 \qquad\text{so}\qquad
 e^y\bigl(z+\sqrt{1+z^2}\bigr)=\frac{1+\sqrt{4D^2+8D+1}}{2D}=\mu .
\]
This is the identity, since
\(\operatorname{arsinh}z=\log(z+\sqrt{1+z^2})\). For the lower bound,
\(\operatorname{arsinh}(\frac12\sinh y)\ge\frac y2\) is equivalent to
\(\frac12\sinh y\ge\sinh\frac y2\), that is, to
\(\sinh\frac y2\cosh\frac y2\ge\sinh\frac y2\). Hence
\(3\log\mu\ge\frac92y=\frac94g_D\).

For the upper bound, \(\frac{D+2}D=(1+\frac1{D+1})/(1-\frac1{D+1})\)
gives \(g_D=2\operatorname{artanh}\frac1{D+1}\ge\frac2{D+1}\). Using
\(\operatorname{arsinh}z\le z\),
\[
 \omega_D=\frac32+\frac{3\operatorname{arsinh}(\frac12\sinh y)}{g_D}
 \le\frac32+\frac{3\sinh y}{2g_D}
 \le\frac32+\frac34\cdot\frac{D+1}{\sqrt{D(D+2)}} .
\]
Finally, \(D(D+2)=(D+1)^2-1\) lies between \(\frac89(D+1)^2\) and
\((D+1)^2\), as \(D\ge2\). Hence
\(\frac{D+1}{\sqrt{D(D+2)}}\le\frac{(D+1)^2}{D(D+2)}
=1+\frac1{D(D+2)}\le1+\frac9{8(D+1)^2}\).
\end{proof}

The expansion
\(\operatorname{arsinh}(\frac12\sinh y)=\frac y2+\frac{y^3}{16}+O(y^5)\)
gives \(\omega_D-\frac94=\frac3{32D^2}+O(D^{-3})\), about a ninth of the
upper bound.

\begin{lemma}[Defect]\label{lem:defect}
Let \(D\ge2\), \(R=8D\) and \(L\ge0\). Starting at \((1,1)\), build
the ramp of Lemma~\ref{lem:ramp} up to \((R,R)\), and then take \(L\)
further smallest parents of arity \(D\). Read forwards, this is a chain
with terminal pair \((1,1)\) and \(J=16D-2+L\) stages: \(L\) bulk
stages of arity \(D\), then the \(2R-2\) ramp stages. All its arities
satisfy \(2\le d_i\le24D-3\). Its gain is \(G=\log F(R,R)+Lg_D\), and
\begin{equation}\label{eq:defect}
 3\log s-\tfrac94G<3\log\tfrac{19}8+Lg_D\bigl(\omega_D-\tfrac94\bigr)
 \le3+\frac{27L}{16D(D+1)^2}.
\end{equation}
\end{lemma}

We call this chain the \emph{ramp-and-bulk chain} for \(D\) and \(L\).

\begin{proof}
The stages, arities and gain follow from
Lemmas~\ref{lem:smallest-parent} and~\ref{lem:ramp}. Since
\(\frac94R+D=\frac{19}8R\), Lemma~\ref{lem:bulk-potential}(b) gives
\(3\log s<3\log\frac{19R}8+3L\log\mu\), and Lemma~\ref{lem:ramp} gives
\(\frac94\log F(R,R)\ge3\log R\). Subtracting proves the first
inequality, since \(3\log\mu-\frac94g_D=g_D(\omega_D-\frac94)\). For the
second, \(\log\frac{19}8<1\), and
\(g_D(\omega_D-\frac94)\le\frac2D\cdot\frac{27}{32(D+1)^2}\) by
Lemma~\ref{lem:ideal-gap} and \(g_D=\log(1+\frac2D)\le\frac2D\).
\end{proof}

\subsection{Types, and the least profile as a potential}
\label{subsec:chains-potential}

For \(d\ge2\) and \(N\ge1\), a type \((\nK,\nO,\nmid)\) of width \(N\),
that is, with \(\nK+\nO+\nmid=N\), has the sector count
\(N!\,(d-1)^{\nK}2^{\nO}/(\nK!\,\nO!\,\nmid!)\)
(Lemma~\ref{lem:stage-types}). Let \(M(N,d)\) be the largest sector
count over all types of width \(N\).

\begin{lemma}[Types]\label{lem:types}
For \(d\ge2\) and \(N\ge1\), every sector count is less than
\((d+2)^N\), and \((d+2)^N/(N+1)^2\le M(N,d)<(d+2)^N\).
\end{lemma}

\begin{proof}
By the multinomial theorem, the sector counts of all types add up to
\(((d-1)+2+1)^N=(d+2)^N\). There are \(\binom{N+2}2\ge3\) types, and
every count is positive, so each count is less than the sum. The
largest count is at least the average, and
\(\binom{N+2}2=\frac{(N+2)(N+1)}2\le(N+1)^2\).
\end{proof}

\begin{lemma}[The least profile is a potential]
\label{lem:profile-potential}
Let \(F\) be the least profile of~\citep[(3.4)]{FinCon26}.
\begin{enumerate}
\item[(a)] Every legal stage \((r,b)\to(c,r)\) of arity \(d\) satisfies
\(d\,F(r,b)\ge(d+2)\,F(c,r)\), with equality on every stage of the ramp
of Lemma~\ref{lem:ramp}.
\item[(b)] Every chain (Section~\ref{subsec:chain-counts}) satisfies
\(\log t+G\le\log F(r_1,b_1)\le\frac43\log s\).
\item[(c)] Every chain scheme of Section~\ref{subsec:chains-legal}, and
every relative scheme given by the variant of
Remark~\ref{rem:varying-powers}, satisfies
\(p>3\log s/(\log t+G)\ge9/4\).
\end{enumerate}
\end{lemma}

\begin{proof}
(a) Put \(b^+=db-(d-1)c\) and \(h=b-c\), so that \(c<b<b^+\) and
\(b^+-b=(d-1)h\). By~\citep[Theorem~3.2]{FinCon26}, \(F\) is a profile:
it satisfies the concavity axiom \citep[(3.2)]{FinCon26}, and the
tripling axiom \citep[(3.3)]{FinCon26} at every column. By
concavity, the increments \(F(r,b'+1)-F(r,b')\) of row \(r\) do not
increase in \(b'\). So each of the \(h\) increments from column \(c\) to
\(b\) is at least each of the \((d-1)h\) increments from \(b\) to
\(b^+\), and \((d-1)\bigl(F(r,b)-F(r,c)\bigr)\ge F(r,b^+)-F(r,b)\), that
is,
\[
 d\,F(r,b)\ge(d-1)F(r,c)+F(r,b^+).
\]
By \citep[(3.5)]{FinCon26}, the increments of row \(r\) from column
\(r-1\) on equal \(\varsigma_r>0\). As they do not increase, the row is
increasing. By Lemma~\ref{lem:reset}(c), \(b^+\ge3c+r-1\), so
\(F(r,b^+)\ge F(r,3c+r-1)\ge3F(r,c)\) by tripling. Hence
\(dF(r,b)\ge(d+2)F(r,c)\), and \(F(r,c)=F(c,r)\) by symmetry.

On the ramp, \citep[(3.4)]{FinCon26} gives
\(F(j,j)=\varsigma_j\frac{3j-1}2\), \(F(j,j+1)=\varsigma_j\frac{3j+1}2\) and
\(F(j+1,j+1)=\varsigma_{j+1}\frac{3j+2}2=\varsigma_j\frac{(3j+1)(3j+2)}{6j}\).
Hence \((3j-1)F(j,j+1)=(3j+1)F(j,j)\) and
\(3j\,F(j+1,j+1)=(3j+2)F(j,j+1)\): these are the equality cases for the
stages of arities \(3j-1\) and \(3j\).

(b) Applying (a) to every stage and multiplying gives
\(F(r_1,b_1)\ge\exp(G)F(c_J,r_J)=t\exp(G)\), by the normalization
\citep[(3.1)]{FinCon26}. The second inequality is
\(F(m,n)\le(m+n-1)^{4/3}\) \citep[(3.6)]{FinCon26}.

(c) By~\eqref{eq:stage-chain-counts},
\(r_*/m_*=s^{6Nk}\prod_i((2M_i-1)/M_i)^{6k}E_i\), and every factor
other than \(s^{6Nk}\) is at least one, so
\(\log(r_*/m_*)\ge6Nk\log s\). By Lemma~\ref{lem:types},
\(M_i<(d_i+2)^N\), so \(\log(B/A)<2Nk(\log t+G)\), since \(J\ge1\).
As \(B>A\) and \(s\ge b_1>c_1\ge1\), we get
\(p\ge6Nk\log s/\log(B/A)>3\log s/(\log t+G)\). Since \(G>0\), (b)
gives \(3\log s/(\log t+G)\ge9/4\). The same argument, with
\(N_1k_1\) in place of \(Nk\), applies to the variant of
Remark~\ref{rem:varying-powers}, with stage-dependent \(N_i\) and
\(k_i\) and \(N_ik_i=N_1k_1\), since there
\(M_i^{2k_i}<(d_i+2)^{2N_ik_i}=(d_i+2)^{2N_1k_1}\).
\end{proof}

Lemma~\ref{lem:profile-potential} is the analogue for chain schemes of
the scalar limit in~\citep[Remark~3.3]{FinCon26}.
Section~\ref{sec:limits} shows more: \(p>9/4\) persists if the auxiliary
of each stage is replaced by a compressed auxiliary of any rank, and for
every supply (Definition~\ref{def:compressed},
Corollary~\ref{cor:compression-limit}).
On a ramp stage, \(c\ge r-1\), \(b^+=3c+r-1\), and row \(r\) of \(F\) is
affine from column \(r-1\) on, so every inequality in the proof of (a)
is an equality.

\subsection{The main theorem}
\label{subsec:chains-main}

\begin{theorem}[Main theorem]\label{thm:main}
Let \(0<\eps\le1\) be rational. Put
\[
 D=\bigl\lceil\sqrt{4/\eps}\,\bigr\rceil,\qquad R=8D,\qquad
 L=\lceil64D/\eps\rceil,
\]
let \(H\) be the number of binary digits of \(\lceil128D/\eps\rceil\),
and put \(N=k=\lceil256DH/\eps\rceil\). Take the ramp-and-bulk chain
for \(D\) and \(L\) (Lemma~\ref{lem:defect}), the width \(N\) and the
power \(k\), and at stage \(i\) a type of width \(N\) with the largest
sector count \(M_i=M(N,d_i)\). Then the restriction
\(r_*\MM_A\restr m_*\MM_B\) of Proposition~\ref{prop:chain-counts} for
these data has rational coefficients and \(B>A\), so it is a relative
scheme over \(\Q\). Its pooled exponent \(p\), its fixed-batch bound
\(\Phi\) and its counts satisfy
\begin{gather*}
 \frac94<p<\Phi<\frac94+\frac{719}{1988}\,\eps<\frac94+\frac\eps2,
 \qquad
 \log A<3\cdot10^9\,\eps^{-9/2}\log^3\frac{4096}\eps,\\
 B\le A^2,\qquad m_*=B^3\le A^6,\qquad r_*<A^{35/4}.
\end{gather*}
\end{theorem}

All parameters are computed from \(\eps\) by integer arithmetic; for
example, if \(\eps=u/v\) with positive integers \(u\) and \(v\), then
\(D\) is the least integer with \(D^2u\ge4v\).

\begin{proof}
Here \(t=1\), so by~\eqref{eq:p-normalized}
and~\eqref{eq:chains-phi},
\(\Phi=(3\log s+C_{\mathrm{sep}}+C_{\mathrm{ex}}+C_{\mathrm{pool}})/\Delta\)
with \(\Delta=\sum_i(\frac1N\log M_i-\log d_i)\), once we know
\(\Delta>0\). Put \(x=128D/\eps\), \(a_N=2\log(N+1)/N\) and
\(\Lambda=\log(4096/\eps)\).

\emph{Step 0: parameter facts.} Since \(\eps\le1\), we have \(D\ge2\),
\(D^2\eps\ge4\) and \(x\ge256\). From
\(2^{H-1}\le\lceil x\rceil<2^H\) we get \(H>\log_2x\ge\log x\),
\(H\le1+\log_2(x+1)\le x\) and, as \(\lceil x\rceil\ge2^8\), \(H\ge9\).
Next, \(N\ge2xH\), so \(N=k\ge4608\), and
\(N+1\le2xH+2\le2x^2+2\le x^3\). Also \(\log(6k)<k\). Since
\(L\ge64D\), \(J=16D-2+L\le\frac54L\). Every arity satisfies
\(d_i\le24D-3\), so \(\log d_i<\log(d_i+2)<\log(24D)\le\log x<H\).

\emph{Step 1: the denominator, and \(B>A\).} By
Lemma~\ref{lem:types}, \(\frac1N\log M_i\ge\log(d_i+2)-a_N\), so
\(\Delta\ge G-Ja_N\). By Step~0, \(a_N\le6\log x/N<6H/(2xH)=3/x\), that
is, \(a_N<3\eps/(128D)\). Since \(\log(1+u)\ge u/(1+u)\), we have
\(g_D\ge2/(D+2)\). As \(\log F(R,R)\ge0\),
\[
 G\ge Lg_D\ge\frac{2L}{D+2}\ge\frac{128D}{(D+2)\eps}\ge\frac{64}\eps,
 \qquad
 \frac JG\le\frac{5L}4\cdot\frac{D+2}{2L}\le\frac{5D}4 .
\]
Hence \(Ja_N/G\le\frac{15}{512}\eps<1\), and \(\Delta\ge G-Ja_N>0\). So
\(\log(B/A)=2Nk\Delta>0\), that is, \(B>A\).

\emph{Step 2: the charges.} Since \(2-1/M_i<2\),
\(C_{\mathrm{sep}}\le3J\log2/N\). Since \(E_i\le6kM_i\le6k(d_i+2)^N\),
\(C_{\mathrm{ex}}\le J\log(6k)/(2Nk)+\frac1{2k}\sum_i\log(d_i+2)\).

\emph{Step 3: reduction.} Put
\(X=3\log s+C_{\mathrm{sep}}+C_{\mathrm{ex}}+C_{\mathrm{pool}}-\frac94\Delta\),
so that \(\Phi-\frac94=X/\Delta\). By
Lemma~\ref{lem:profile-potential}(c), \(p>9/4\), that is,
\(X>C_{\mathrm{pool}}>0\). With \(\Delta\ge G-Ja_N\) and
\(-\frac94\Delta\le-\frac94G+\frac94Ja_N\) from Step~1,
\[
 \Phi-\frac94=\frac X\Delta\le\frac X{G-Ja_N}
 \le\frac1{1-Ja_N/G}\Bigl(\frac{3\log s-\frac94G}G
 +\frac{C_{\mathrm{sep}}+C_{\mathrm{ex}}+C_{\mathrm{pool}}}G
 +\frac94\cdot\frac{Ja_N}G\Bigr).
\]

\emph{Step 4: the charges, divided by \(G\).} We bound the six terms of
the bracket: (i) the geometric defect, (ii) the type loss, (iii) the
separator copies, (iv) and (v) the two parts of the exactification
charge, and (vi) the fixed-batch ceiling.
\begin{enumerate}
\item[(i)] By~\eqref{eq:defect},
\(3\log s-\frac94G<3+\frac{27L}{16D(D+1)^2}\). Here
\(3/G\le\frac3{64}\eps\), and \(G\ge2L/(D+2)\) gives
\(\frac{27L}{16D(D+1)^2G}\le\frac{27(D+2)}{32D(D+1)^2}<\frac{27}{32D^2}
\le\frac{27}{128}\eps\), using \(D(D+2)<(D+1)^2\) and
\(D^2\eps\ge4\). Together, \(\frac3{64}+\frac{27}{128}=\frac{33}{128}\).
\item[(ii)] \(\frac94\cdot\frac{Ja_N}G\le\frac94\cdot\frac{15}{512}\eps\),
by Step~1.
\item[(iii)] \(\frac{3J\log2}{NG}\le\frac{15D\log2}{4N}
\le\frac{15\log2}{1024H}\eps<\frac{15}{1024}\eps\), by \(J/G\le5D/4\) and
\(N\ge256DH/\eps\).
\item[(iv)] \(\frac{J\log(6k)}{2NkG}\le\frac J{2NG}\le\frac{5D}{8N}
\le\frac5{2048H}\eps\le\frac5{2048}\eps\).
\item[(v)] \(\frac1{2kG}\sum_i\log(d_i+2)\le\frac{JH}{2kG}
\le\frac{5DH}{8k}\le\frac5{2048}\eps\).
\item[(vi)] \(\frac{C_{\mathrm{pool}}}G=\frac{\log2}{2NkG}
\le\frac{\log2}{128Nk}\eps<\frac1{128}\eps\), by \(G\ge64/\eps\).
\end{enumerate}
The bounds add up to
\(\frac{528+135+30+5+5+16}{2048}\eps=\frac{719}{2048}\eps\), and
\(1-Ja_N/G\ge\frac{497}{512}\). Hence
\(\Phi-\frac94<\frac{719}{2048}\cdot\frac{512}{497}\eps
=\frac{719}{1988}\eps<\frac\eps2\). With \(9/4<p\)
(Lemma~\ref{lem:profile-potential}(c)) and
\(\Phi=p+C_{\mathrm{pool}}/\Delta>p\), this proves the exponent bounds.

\emph{Step 5: the source side.} We have
\(D\le2\eps^{-1/2}+1\le3\eps^{-1/2}\), so
\(J\le16D+64D/\eps-1\le240\,\eps^{-3/2}\) and
\(x\le384\,\eps^{-3/2}\). Since
\(H\le1+\log_2(2x)\le1+\frac32\log(2x)\) and \(\log(2x)\ge\log512>2\),
we get \(H\le2\log(2x)\le2\log768+3\log\frac1\eps\le3\Lambda\), using
\(768^2<4096^3\). Hence
\(N=k\le256DH/\eps+1\le2304\,\eps^{-3/2}\Lambda+1
\le2305\,\eps^{-3/2}\Lambda\). Finally,
\(\log d_i<\log(24D)\le\log(72\eps^{-1/2})\le\Lambda\). Therefore
\[
 \log A=2Nk\sum_i\log d_i<2N^2J\Lambda
 \le2\cdot2305^2\cdot240\,\eps^{-9/2}\Lambda^3
 =2{,}550{,}252{,}000\,\eps^{-9/2}\Lambda^3
 <3\cdot10^9\,\eps^{-9/2}\Lambda^3 .
\]

\emph{Step 6: the other counts.} Since \(d_i\ge2\), we have
\(d_i+2\le d_i^2\), so \(M_i<(d_i+2)^N\le d_i^{2N}\) by
Lemma~\ref{lem:types}, and \(B=\prod_iM_i^{2k}\le A^2\). Since \(t=1\),
\eqref{eq:stage-chain-counts} gives \(m_*=B^3\le A^6\). By the definition
of \(p\), \(r_*=m_*(B/A)^p\). As \(1<B/A\le A\) and
\(p<\frac94+\frac\eps2\le\frac{11}4\), we get
\(r_*<B^3(B/A)^{11/4}\le A^{35/4}\). The coefficients are rational by
Proposition~\ref{prop:chain-counts}.
\end{proof}

With the separator of~\citep[Proposition~4.6]{FinCon26} in place of
Lü's, \(E_i=1+6k(M_i-1)^2\) doubles bound~(v) of Step~4, so the proof
gives \(\frac{724}{1988}\eps\) in place of \(\frac{719}{1988}\eps\),
with coefficients in a cyclotomic field instead of~\(\Q\). For the bulk
alone, started at \((1,1)\), Lemma~\ref{lem:bulk-potential}(b) with
\(R=1\) puts \(3\log(D+\frac94)\) in place of \(3\log\frac{19}8\)
in~\eqref{eq:defect}; the same argument would then need \(L\) larger by
a factor of order \(\log\frac1\eps\) and would give only
\(\log A=O(\eps^{-9/2}\log^4\frac1\eps)\).

\begin{corollary}\label{cor:main-eps}
Let \(0<\eps\le\frac12\) be rational. The scheme of
Theorem~\ref{thm:main} with \(2\eps\) in place of \(\eps\) is a
relative scheme over \(\Q\) with \(\frac94<p<\Phi<\frac94+\eps\),
\(\log A<1.4\cdot10^8\,\eps^{-9/2}\log^3(2048/\eps)\), \(B\le A^2\),
\(m_*\le A^6\) and \(r_*<A^{35/4}\).
\end{corollary}

\begin{proof}
Theorem~\ref{thm:main} applies to \(2\eps\le1\). It gives
\(\Phi<\frac94+\frac{719}{994}\eps\) and
\(\log A<3\cdot10^9(2\eps)^{-9/2}\log^3(2048/\eps)\). Finally,
\(3\cdot10^9\cdot2^{-9/2}<1.4\cdot10^8\), because
\(9\cdot10^{18}<1.96\cdot10^{16}\cdot2^9\).
\end{proof}

\begin{remark}[\(\Phi\) versus \(p\), and the bound \(\frac94+\frac\eps2\)]
\label{rem:phi-p}
By~\eqref{eq:prelim-phi}, \(\Phi\) also bounds the exponent of the
recursion without pooling. Here \(\Phi-p=\log2/\log(B/A)\), and Steps~0
and~1 give
\(\log(B/A)\ge2Nk(G-Ja_N)\ge2\cdot4608^2\cdot\frac{497}{512}\cdot
\frac{64}\eps\), so \(\Phi-p<10^{-9}\eps\). The theorem keeps half of
\(\eps\) in reserve because Section~\ref{sec:seeds} unrolls the pooled
recursion into one seed, as in~\citep[Theorem~9.2]{FinCon26}; the factor
lost in unrolling is absorbed into the remaining \(\eps/2\). For an
algorithm with exponent below \(\frac94+\eps\) alone,
Corollary~\ref{cor:main-eps} suffices.
\end{remark}
